Un sistema adimensional satisface
\[
y''+3y'+2y=u(t),\qquad y(0)=y'(0)=0.
\]
Se sabe que la entrada es un único pulso de altura finita:
\[
u(t)=A\quad(0\le t<\tau),\qquad
u(t)=0\quad(t>\tau),\qquad A>0,\quad0<\tau<1.
\]
Se desconoce tanto \(A\) como \(\tau\). El registro sintético de la respuesta sólo está disponible después de \(t=1\) y se trata como exacto:
\[
y(t)=2e^{-t}-3e^{-2t}\qquad(t\ge1).
\]
Recupere la altura y la duración del pulso, y reconstruya la respuesta anterior al registro. Use explícitamente las condiciones de unión que impone una entrada finita en \(t=\tau\); indique qué derivadas pueden cambiar de valor al conmutar. El valor aislado que se asigne a la entrada en el instante del cambio no forma parte de la inferencia.
